\section{Identity systems and the original flattening construction}
\label{sec:p07-hott-identity-systems-and-flattening}

This development follows the June 2013 public-release
\cite[\S5.8, \S6.12, \S8.1]{HoTT2013PublicReleaseFoundations}.
It distinguishes identity-system induction from identity elimination in a
fixed loop type, and proves flattening for the source's homotopy-coequalizer
presentation, not for an unspecified class of all higher inductive types.
Use the routed path, transport, Sigma-path, funext, univalence and circle
interfaces, including the ordinary integer presentation in
Section~\ref{sec:p07-hott-circle-encode-decode}.
Composition \(p\cdot q\) traverses \(p\) first. All path comparisons below
are propositional unless explicitly marked \(\equiv_{\mathrm j}\).

\begin{hbdefinition}{def:p07-is-pointed}{Pointed predicates and identity-system induction}
For \(a_0:A\), a pointed predicate is \(R:A\to\mathcal U\) with
\(r_0:R(a_0)\). For \((S,s_0)\), put
\[
\mathsf{ppmap}(R,S)=
 \sum_{u:\prod_{b:A}(R(b)\to S(b))}(u(a_0,r_0)=s_0).
\]
An identity system at \(a_0\) supplies, for every
\(D:\prod_{b:A}(R(b)\to\mathcal V)\) and \(d:D(a_0,r_0)\),
a section \(v:\prod_{b:A}\prod_{r:R(b)}D(b,r)\) and a path
\(v(a_0,r_0)=d\). Existence here is supplied data, not propositional
truncation of a section. The base computation is propositional, unlike
the primitive judgmental base computation of \(J\).
\end{hbdefinition}

\begin{hbtheorem}{thm:p07-is-characterization}{proved=annex:p07-is-characterization}{Four characterizations, with the pointed computation retained}
The following conditions on \((R,r_0)\) imply one another:
identity-system induction; contractibility of every
\(\mathsf{ppmap}(R,S)\); equivalence of every canonical map
\[
e_b:(a_0=b)\to R(b),\qquad e_b(p)=T_R(p,r_0);
\]
and contractibility of \(\sum_{b:A}R(b)\).
The proof constructs the pointed-map equality, inverse homotopies, total
space comparison and normalized base computation.
Funext is used for function paths. No setness of \(A\) is assumed.
\end{hbtheorem}

\begin{hbdefinition}{def:p07-is-two-sided}{Two-sided identity systems}
A two-sided identity system is \(R:A\to A\to\mathcal U\) with
\(r_0:\prod_aR(a,a)\), supplying sections for every
\(D:\prod_{a,b:A}(R(a,b)\to\mathcal V)\), given diagonal data
\(d:\prod_aD(a,a,r_0(a))\), with propositional diagonal computation.
In this definition the bound \(r\) has type \(R(a,b)\).
\end{hbdefinition}

\begin{hbtheorem}{thm:p07-is-two-sided}{proved=annex:p07-is-two-sided}{Two-sided induction, local systems and coherent maps}
Two-sided induction is equivalent to identity-system induction for every
\((R(a,-),r_0(a))\). These conditions are also equivalent to:
contractibility, for every \((S,s_0)\), of
\[
\sum_{u:\prod_{a,b:A}(R(a,b)\to S(a,b))}
       \prod_a(u(a,a,r_0(a))=s_0(a));
\]
equivalence of every \(p\mapsto T_{R(a,-)}(p,r_0(a))\);
and contractibility of every \(\sum_bR(a,b)\).
\end{hbtheorem}

\begin{hbcorollary}{thm:p07-is-equivalence-homotopy-induction}{proved=annex:p07-is-equivalence-homotopy-induction}{Equivalence and homotopy induction}
Univalence supplies two-sided identity-system induction on
\(R(X,Y)=(X\simeq Y)\), with diagonal identity equivalence.
Funext supplies it on \(R(f,g)=\prod_x(f(x)=g(x))\), with diagonal
pointwise reflexivity. Thus a family over equivalences, respectively
homotopies, is eliminated from its identity, respectively reflexive,
case with a propositional computation path. Universe levels must allow
the stated motive; no resizing is inferred.
\end{hbcorollary}

\begin{hbdefinition}{def:p07-flat-coequalizer}{The exact coequalizer and flattened HIT principles}
Let \(f,g:B\to A\). The source's \(W\) has
\(\mathsf c:A\to W\) and
\(\mathsf p:\prod_{b:B}(\mathsf c(fb)=\mathsf c(gb))\).
For \(F:W\to\mathcal V\), its dependent eliminator takes
\(u:\prod_aF(\mathsf c(a))\) and
\(\prod_b(T_F(\mathsf p_b,u(fb))=u(gb))\), returning a section
with judgmental point and propositional dependent-path computation.
These HIT principles are additional type formers, not MLTT existence
or normalization results.

Let \(C:A\to\mathcal U\), and
\(D:\prod_b(C(fb)\simeq C(gb))\).
Here \(D\) is a family of equivalences, not the motive named \(D\)
in identity-system induction. Recursion into a sufficiently large
universe gives \(P:W\to\mathcal U\) with
\[
P(\mathsf c(a))\equiv_{\mathrm j}C(a),\qquad
\operatorname{ap}_P(\mathsf p_b)=\operatorname{ua}(D_b).
\]
The latter computation is propositional.
The flattened type \(\widetilde W\) has constructors
\[
\widetilde{\mathsf c}:\prod_a(C(a)\to\widetilde W),\qquad
\widetilde{\mathsf p}:\prod_b\prod_{y:C(fb)}
 \bigl(\widetilde{\mathsf c}(fb,y)
       =\widetilde{\mathsf c}(gb,D_b(y))\bigr).
\]
Its dependent eliminator takes point data at
\(\widetilde{\mathsf c}(a,y)\) and a transported path for each
\(\widetilde{\mathsf p}_{b,y}\), with the same point/path
computation standing. No higher constructor is imposed here.
\end{hbdefinition}

\begin{hbconvention}{def:p07-flat-pairpaths}{Pair paths and calibrated comparisons}
Here \(Y:X\to\mathcal U\), and
\(Q:(\sum_{x:X}Y(x))\to\mathcal V\);
write \(Q(x,u)=Q((x,u))\).
Write \([s,r]=\operatorname{pairpath}(s,r):(x,u)=(x',u')\)
when \(s:x=x'\) and \(r:T_Y(s,u)=u'\).
Let \(H(s,u)=[s,\mathsf{refl}_{T_Y(s,u)}]\) and
\(V_x(r)=[\mathsf{refl}_x,r]\).
For \(F(x)=\prod_{u:Y(x)}Q(x,u)\), denote by \(\mathcal K_s\) the
dependent-function-path equivalence
\[
(T_F(s,v)=v')\simeq
\prod_u\bigl(T_Q(H(s,u),v(u))=v'(T_Y(s,u))\bigr).
\]
It is chosen by \(J\) on \(s\), reducing at reflexivity to the actual
happly/funext equivalence. Its two inverse comparisons are retained.
For constant \(Q=Z\), denote its normalized specialization by
\(\mathcal K^0_s\), with target
\(\prod_u(v(u)=v'(T_Y(s,u)))\).
\end{hbconvention}

\begin{hblemma}{thm:p07-flat-pair-calculus}{proved=annex:p07-flat-pair-calculus}{Pair composition, two-paths and action of a curried function}
The calibrated comparisons above have both inverse homotopies.
Pair paths obey
\[
[s,r]\cdot[t,v]
 =[s\cdot t,\,
   \gamma_Y(s,t,u)\cdot
   \operatorname{ap}_{T_Y(t,-)}(r)\cdot v],
\qquad [s,r]=H(s,u)\cdot V_{x'}(r).
\]
Here \(\gamma_Y(s,t,u):T_Y(s\cdot t,u)=T_Y(t,T_Y(s,u))\).
For \(d:\prod_{x:X}(Y(x)\to Z)\) into a constant type \(Z\),
put \(L(x,u)=d(x)(u)\). Then
\[
\operatorname{ap}_L([s,r])
 =\mathcal K^0_s(\operatorname{apd}_d(s))(u)
       \cdot\operatorname{ap}_{d(x')}(r).
\]
Changing the base-path coordinate from \(s\) to \(s'\) by
\(\alpha:s=s'\) requires transporting the second coordinate in
\(s\mapsto(T_Y(s,u)=u')\); it cannot be ignored.
\end{hblemma}

\begin{hblemma}{thm:p07-flat-total-constructors}{proved=annex:p07-flat-total-constructors}{The actual total-space generators}
Put \(T=\sum_{w:W}P(w)\). There are
\[
\mathsf c'(a,y)=(\mathsf c(a),y),\qquad
\mathsf p'_{b,y}=[\mathsf p_b,\theta_b(y)],
\quad
\theta_b(y):T_P(\mathsf p_b,y)=D_b(y).
\]
The path \(\theta_b(y)\) is constructed from universe recursion and
univalence's transport comparison, not judgmental computation.
\end{hblemma}

\begin{hblemma}{thm:p07-flat-total-induction}{proved=annex:p07-flat-total-induction}{Dependent total-space induction}
For \(Q:T\to\mathcal V\), suppose
\(a_{a,y}:Q(\mathsf c'(a,y))\) and
\[
b_{b,y}:T_Q(\mathsf p'_{b,y},a_{fb,y})
                       =a_{gb,D_b(y)}.
\]
There is a section \(v:\prod_{z:T}Q(z)\) with
\(v(\mathsf c'(a,y))\equiv_{\mathrm j}a_{a,y}\).
The proof explicitly constructs the dependent-function lift on every
\(\mathsf p_b\). This statement does not assert an additional general
judgmental or propositional computation of
\(\operatorname{apd}_v(\mathsf p'_{b,y})\).
\end{hblemma}

\begin{hblemma}{thm:p07-flat-total-recursion}{proved=annex:p07-flat-total-recursion}{Recursion with the necessary path computation}
For constant \(Z\), point data \(a_{a,y}:Z\), and edge data
\(b_{b,y}:a_{fb,y}=a_{gb,D_b(y)}\), there is \(v:T\to Z\) with
judgmental point computation and
\(\operatorname{ap}_v(\mathsf p'_{b,y})=b_{b,y}\).
This latter path computation is proved by the curried-action formula
and cancellation of the two explicitly identified fiber paths.
\end{hblemma}

\begin{hbtheorem}{thm:p07-flat-equivalence}{proved=annex:p07-flat-equivalence}{The original coequalizer flattening lemma}
There is an equivalence
\[
\left(\sum_{w:W}P(w)\right)\simeq\widetilde W.
\]
Its two maps are recursion on \(\widetilde W\) using
\(\mathsf c',\mathsf p'\), and total-space recursion using
\(\widetilde{\mathsf c},\widetilde{\mathsf p}\).
Both inverse homotopies have their dependent edge lifts supplied.
\end{hbtheorem}

\begin{hbdefinition}{def:p07-flat-circle-cover}{The native integer-indexed cover}
For \(A=B=1\), \(f=g=\mathrm{id}\), the coequalizer presentation
is equivalent to the existing circle, not asserted judgmentally equal
to it. For \(C=\mathbb Z\) and \(D=e_{\mathsf{succ}}\), the
single-point specialization of the flattening construction, denoted
\(R\) in the source's circle application,
has \(\mathsf c:\mathbb Z\to R\) and
\(\mathsf p_n:\mathsf c(n)=\mathsf c(\mathsf{succ}(n))\).
This \(R\) is a HIT, not the pointed predicate of the first subsection;
the overlap is the native source notation and is locally typed.
\end{hbdefinition}

\begin{hbtheorem}{thm:p07-flat-circle-cover}{proved=annex:p07-flat-circle-cover}{Flattened cover and its signed contraction}
The circle-specialized flattening construction gives
\(\sum_{x:S^1}\mathrm{code}(x)\simeq R\).
The type \(R\) is contractible at \(\mathsf c(0)\); its contraction
uses integer-indexed paths \(q_n:\mathsf c(0)=\mathsf c(n)\)
and proves \(q_n\cdot\mathsf p_n=q_{\mathsf{succ}(n)}\) in every
signed case. Consequently the total code space is contractible by
this route independently of either circle inverse homotopy.
\end{hbtheorem}

\begin{hblemma}{thm:p07-is-circle-pair-coherence}{proved=annex:p07-is-circle-pair-coherence}{The source's omitted winding pair-path coherence}
Let \(h_n:T_{\mathrm{code}}(\mathsf{loop}^n,0)=n\) be the
already proved signed winding path and let
\(\theta_n:T_{\mathrm{code}}(\mathsf{loop},n)=\mathsf{succ}(n)\).
Put \(K_n=[\mathsf{loop}^n,h_n]\) and
\(S_n=[\mathsf{loop},\theta_n]\). Then
\[
K_0=\mathsf{refl}_{(\mathsf{base},0)},\qquad
K_n\cdot S_n=K_{\mathsf{succ}(n)}.
\]
The proof accounts for the first-coordinate power comparison and the
transported second-coordinate winding proof. Integer setness, not
setness of circle loops, provides the necessary fiber two-path.
\end{hblemma}

\begin{hbtheorem}{thm:p07-is-circle-native-induction}{proved=annex:p07-is-circle-native-induction}{The native identity-system route for the code}
The pointed predicate \((\mathrm{code},0)\) is an identity system at
\(\mathsf{base}\). Given \(D\) and \(d:D(\mathsf{base},0)\), its
base-fiber datum is exactly
\[
a(n)=T_{\widehat D}(K_n,d),
\qquad
\widehat D:\left(\sum_x\mathrm{code}(x)\right)\to\mathcal V.
\]
The preceding pair coherence constructs its circle lift, and \(K_0\)
gives the pointed computation. The general characterization therefore
identifies the canonical transport encoder as an equivalence, without
rerunning its separate encode--decode inverse proof.
\end{hbtheorem}

\begin{hbexample}{def:p07-is-failure-boundaries}{A pointed family need not be an identity system}
The constant pointed predicate \(R(x)=1\) over \(S^1\), pointed by
\(\star\), does not satisfy identity-system induction.
Its total space projects equivalently to \(S^1\): the inverse is
\(x\mapsto(x,\star)\), and unit induction supplies the second
inverse homotopy. If that total space were contractible, so would
the circle be; its loop type would then be contractible, contrary
to the already established \(\Omega(S^1)\simeq\mathbb Z\) and the
disjoint integer tags \(0,1\). Equivalently the canonical map
\(\Omega(S^1)\to1\) cannot be an equivalence.
Flattening this constant family produces a point with a loop, not a
contractible interval. The two inverse maps in the general flattening
theorem must retain every edge constructor; point agreement alone
does not supply the needed homotopy lift. The integer-successor
family's signed contraction is a special proved property.
\end{hbexample}

\begin{proofboundary}{def:p07-is-flattening-boundary}{Exact theorem and computation scope}
The coequalizer flattening theorem has no recursive point constructor
or higher path constructor in its hypotheses. Its proof does not assert
a general syntactic admissibility theorem, a normalization/canonicity
result, or flattening for every possible HIT. The dependent total-space
induction theorem asserts its point computation; the required
nondependent edge computation is separately proved. None is silently
made judgmental. The two-sided induction corollaries use the specified
funext or univalence axiom and their universe levels, not UIP, choice
or resizing. Computational cubical composition has the distinct
presentation of Section~\ref{sec:p07-cchm2016}. Higher homotopy-group
calculations additionally require their truncation and iterated-loop
interfaces; they do not follow from this coequalizer flattening theorem.
\end{proofboundary}
